
\UserCommand{Add Section Link}<\cs{seclink}>{
    \cs{seclink}
    \{\meta{section path}\}
    (\meta{link text})
}[
    - \meta{section path} : The path of the target section.
][
    - \meta{link text} : The text to be displayed for the link.
]

\UserCommand{Add Block Link}<\cs{blclink}>{
    \cs{blclink}
    \{\meta{label}\}
    (\meta{link text})
}[
    - \meta{label} : The label of the target block.
][
    - \meta{link text} : The text to be displayed for the link.
]

\UserCommand{Add Block Index Link}<\cs{indlink}>{
    \cs{indlink}
    \{\meta{scope path}\}
    (\meta{link text})
}[
    - \meta{scope path} : The scope path of the target block index.
][
    - \meta{link text} : The text to be displayed for the link.
]

The link will be rendered using \meta{link text} as its display content. If this argument is omitted, \SimpleSystemTeX will apply a default link display style. You can modify the relevant display logic by altering the style code in the \seclink{Preamble/StyleConfiguration}(Style Configuration).

\newpage
